% !TeX root = ../../Book3.tex
% 纲要标题：### A.3 素数处的局部化与完备化

\section{素数处的局部化与完备化}
\label{b3:sec:A03}

% 纲要细目（与计划纲要同步）：
% * 局部化和完备化在素数处的作用
%
% <!-- 短节：不设小节 -->


固定数域整数环 \(R=\mathcal O_K\) 中的非零素理想 \(\mathfrak p\)。局部化 \(R_{\mathfrak p}\) 允许把所有不属于 \(\mathfrak p\) 的元素放到分母中。因此它保留的是在这个素理想处的整除次数；一个元素在别的素理想处是否有分母，不再妨碍它属于 \(R_{\mathfrak p}\)。

\begin{definition}\label{b3:def:number-field-prime-completion}
设 \(K\) 为数域，\(\mathcal O_K\) 为其代数整数环，\(\mathfrak p\subset\mathcal O_K\) 为非零素理想。记 \(v_{\mathfrak p}:K^\times\to\mathbb Z\) 为 DVR \((\mathcal O_K)_{\mathfrak p}\) 的正规化赋值，即其参数的赋值为 \(1\)。记
\[
\widehat{\mathcal O}_{K,\mathfrak p}
=\varprojlim_r(\mathcal O_K)_{\mathfrak p}/
(\mathfrak p(\mathcal O_K)_{\mathfrak p})^r,
\qquad
K_{\mathfrak p}=\operatorname{Frac}(\widehat{\mathcal O}_{K,\mathfrak p}).
\]
分别称它们为相应的完备整数环与 \(K\) 在 \(\mathfrak p\) 处的完备化。
\end{definition}

由 DVR 的刻画和正则局部环的完备化定理，\(\widehat{\mathcal O}_{K,\mathfrak p}\) 仍为 DVR，原来的参数仍为参数，剩余域不变。原环嵌入完备化是 Krull 交定理的结果。赋值也因而延拓到 \(K_{\mathfrak p}\)。

\ProseParagraphBreak
局部化和完备化不是同一操作。例如 \(\mathbb Z_{(p)}\) 只由分母不被 \(p\) 整除的有理数组成，是可数集；\(\mathbb Z_p\) 则包含任意相容的 \(p\) 进数位串，是不可数集。二者有相同的模 \(p^r\) 商，却并不相等。这正是完备化补入相容极限的含义。

\begin{theorem}\label{b3:thm:number-field-completion-product}
设 \(K\) 为数域，\(R=\mathcal O_K\) 为其代数整数环，\(p\) 为有理素数，且
\(pR=\prod_{i=1}^r\mathfrak p_i^{e_i}\)。则自然映射给出
\[
R\otimes_{\mathbb Z}\mathbb Z_p
\simeq\prod_{i=1}^r\widehat{\mathcal O}_{K,\mathfrak p_i},
\qquad
K\otimes_{\mathbb Q}\mathbb Q_p
\simeq\prod_{i=1}^rK_{\mathfrak p_i}.
\]
若 \(f_i=[R/\mathfrak p_i:\mathbb F_p]\)，则
\([K_{\mathfrak p_i}:\mathbb Q_p]=e_if_i\)。
\end{theorem}
\begin{proof}
\(R\) 是有限自由 \(\mathbb Z\)-模，有限模的完备化与张量积比较给出
\(R\otimes\mathbb Z_p\simeq\varprojlim_a R/p^aR\)。中国剩余定理给出
\[
R/p^aR\simeq\prod_i R/\mathfrak p_i^{ae_i}.
\]
环 \(R/\mathfrak p_i^b\) 只有一个极大理想；不在 \(\mathfrak p_i\) 中的元素在这个商中已经可逆。因此它与
\(R_{\mathfrak p_i}/\mathfrak p_i^bR_{\mathfrak p_i}\) 相同。指数 \(ae_i\) 在所有正整数中共尾，取逆极限即得第一式。

在每个因子中 \(p\) 等于单位乘以参数的 \(e_i\) 次幂，故把 \(p\) 变成可逆元就得到该 DVR 的分式域。对第一式把 \(p\) 变成可逆元，便得第二式。

每个完备整数环是有限自由 \(\mathbb Z_p\)-模：它是有限自由模 \(R\otimes\mathbb Z_p\) 的直和因子，因而有限且无挠，在 PID \(\mathbb Z_p\) 上自由。其秩等于模 \(p\) 商的 \(\mathbb F_p\)-维数；该商是
\(R/\mathfrak p_i^{e_i}\)，维数为 \(e_if_i\)。再取分式域，得到最后的次数公式。
\end{proof}

此定理解释了为什么在一个有理素数处研究数域时，可能同时出现几个局部域。张量积保留 \(p\) 上方的全部素理想；选定 \(\mathfrak p_i\)，才是选定一个因子。

\begin{example}\label{b3:exm:quadratic-five-completions}
继续取 \(K=\mathbb Q(\sqrt5)\)。模 \(11\) 时 \(T^2-T-1\) 的根 \(4,8\) 都是单根，它们在 \(\mathbb Z_{11}\) 中唯一提升。因此
\[
K\otimes_{\mathbb Q}\mathbb Q_{11}\simeq
\mathbb Q_{11}\times\mathbb Q_{11}.
\]
模 \(2\) 时多项式不可约，只有一个完备化，次数为 \(2\)，分歧指数为 \(1\)，剩余次数为 \(2\)。模 \(5\) 时也只有一个完备化，但此时分歧指数为 \(2\)，剩余次数为 \(1\)。在后者中 \(\sqrt5\) 的赋值为 \(1\)，因为
\(2v_{\mathfrak p}(\sqrt5)=v_{\mathfrak p}(5)=2\)，所以 \(\sqrt5\) 是参数。
\end{example}

这里的单根提升可直接使用\cref{b3:thm:hensel-simple-root}。它帮助计算已经存在的完备因子；各因子的数目与次数，则先由理想分解确定。这样就把第10章的赋值、第13章的完备化与第23章的 Hensel 提升用于同一个算例。
